\documentclass[german,10pt]{book}      \input{../preambel}         \begin{document}  \setcounter{chapter}{0}

\chapter{Der Deutsch-Jozsa-Algorithmus und das Penny-Flip-Spiel}
% Kap x
\label{chap_Deutsch}

\info{Thomas Filk}{15.12.2025}%
Der Deutsch-Jozsa-Algorithmus ist ein vergleichsweise einfacher Quantenalgorithmus,
mit dem man in einer Abfrage entscheiden kann, ob eine Bit-wertige Funktion auf $N$-Bit-Folgen
konstant ist (also allen Folgen entweder den konstanten Wert 0 oder allen Folgen den Wert 1 zuweist)
oder balanciert (in diesem Fall ordnet die Funktion der H\"alfte der Folgen den Wert 0 und
der anderen H\"alfte den Wert 1 zu). M\"ochte man dieses Problem mit einem klassischen
Algorithmus entscheiden, ben\"otigt man im ung\"unstigsten Fall f\"ur eine eindeutige Antwort 
$2^{N-1}+1$ Abfragen. Mit einem 
Quantencomputer kann man zun\"achst s\"amtliche $2^N$ Folgen als Superposition in einem
Zustand kodieren und die gesuchte Funktion -- die man auch als Oracle bezeichnet -- so
auf diesen Zustand anwenden, dass man anschlie\ss end mit einer einzigen Abfrage
entscheiden kann, ob eine konstante oder balancierte Funktion vorliegt. 

Benannt ist der Deutsch-Jozsa-Algorithmus nach David Deutsch und Richard Jozsa, die
die wesentlichen Grundlagen dieses Algorithmus' entwickelt haben \cite{Deutsch,DJ}
(siehe auch \cite{Ekert}). Die mit diesem Algorithmus
gel\"oste Aufgabe ist zwar von keiner praktischen Relevanz, aber das Beispiel zeigt, dass
ein Quantencomputer manche Probleme in wesentlich weniger Schritten l\"osen kann als
ein klassischer Computer, im vorliegenden Fall sind es sogar exponentiell weniger Schritte.

In diesem Kurztext wird auch das Penny-Flip-Spiel beschrieben, da es dem Deutsch-Algorithmus
sehr \"ahnlich ist. Man kann diesen Abschnitt aber auch lesen, ohne vorher den Deutsch-Algorithmus
gelesen zu haben.

\section{Konstante und balancierte Funktionen}
 
Jeder Folge $(b_0,b_1,...,b_{N-1})$ ($b_i\in \{0,1\}$) von $N$ Bits kann man eine 
nat\"urliche Zahl $x$ zwischen 0 und 
$2^N-1$ zuordnen, indem man die Bitfolge als bin\"are Darstellung dieser Zahl interpretiert:
\[                   x = \sum_{i=0}^{N-1} b_i 2^i   \, . \]
Betrachten wir nun Bit-wertige Funktionen $f$ auf solchen Folgen, k\"onnen wir ihre Urbildmenge
als $\{0,1\}^N$ oder $[0,2^N-1]$ ansehen:
\[          f: \{0,1\}^N \rightarrow \{0,1\}  \hspace{1cm} \mbox{bzw.} \hspace{1cm}  f: [0,2^N-1] \rightarrow \{0,1\} \, . \]
$\{0,1\}^N$ bezeichnet das $N$-fache kartesische Produkt der Menge $\{0,1\}$.

Eine konstante Funktion ordnet allen Elementen dasselbe Element
ihrer Bildmenge zu. Im vorliegenden Fall gibt es nur zwei konstante Funktionen, unabh\"angig
von $N$: die Funktion, die allen Elementen die 0 zuordnet, und die Funktion, die allen
Elementen die 1 zuordnet. Man bezeichnet eine Funktion als balanciert (oder auch
ausgeglichen), wenn jedes
Element der Bildmenge gleich h\"aufig als Funktionswert auftritt. Bei einer Bit-wertigen
Funktion wird also der H\"alfte der F\"alle die 0 zugeordnet und der H\"alfte der F\"alle
die 1. 

Insgesamt gibt es $2^{(2^N)}$ Bit-wertige Funktionen auf der Menge der 
$N$-Bit-Folgen.\footnote{Allgemein
gibt es $|B|^{|\Omega|}$ Funktionen auf einer Urbildmenge $\Omega$ mit $|\Omega|$ Elementen
und einer Bildmenge $B$ mit $|B|$ Elementen} Davon sind zwei Funktionen konstant. Von den
balancierten Funktionen gibt es sehr viele (es sind $2^N!/(2^{N-1}!)^2$); die Anzahl ist 
jedoch nicht relevant. Wichtig ist jedoch,
dass man von den $2^N$ Bit-Folgen mit einem klassischen Algorithmus im ung\"unstigsten Fall
$2^{N-1}+1$ (also eine Folge mehr als die H\"alfte aller Folgen) \"uberpr\"ufen muss, um sicher
unterscheiden zu k\"onnen, ob es sich um eine konstante oder eine balancierte Funktion handelt. 
Im g\"unstigsten Fall reichen zwar auch zwei Auswertungen, doch f\"ur eine eindeutige 
Unterscheidung m\"ussen als Funktion von $N$ exponentiell viele F\"alle untersucht werden.
Da mit dem Deutsch-Jozsa-Algorithmus nur eine einzige Auswertung auf einem 
Quantencomputer notwendig ist, ist der Quantencomputer 
exponentiell schneller als der klassische Computer. 

Im Folgenden betrachten wir zun\"achst den Fall $N=1$, d.h.\ die Bit-wertige Funktion ordnet einem
einzelnen Bit wieder ein Bit zu. Von diesen Funktionen gibt es insgesamt vier:
\begin{eqnarray*}
         f_1(0)=0 \hspace{1cm} f_1(1)=0  \hspace{1cm} \mathrm{kurz} & &   f_1(x)=0 \\
         f_2(0)=1 \hspace{1cm} f_2(1)=1  \hspace{1cm} \mathrm{kurz} & &   f_2(x)=1 \\
         f_3(0)=0 \hspace{1cm} f_3(1)=1  \hspace{1cm} \mathrm{kurz} & &   f_3(x)=x \\
         f_4(0)=1 \hspace{1cm} f_1(1)=0  \hspace{1cm} \mathrm{kurz} & &   f_4(x)=1-x =\neg x=x\oplus 1
\end{eqnarray*}
Die ersten beiden Funktionen sind die konstanten Funktionen, die anderen beiden
Funktionen sind balanciert. Die Aufgabe besteht darin, mit m\"oglichst wenigen Abfragen 
an eine unbekannte Funktion -- dem Oracle -- zu entscheiden, ob es sich
bei der unbekannten Funktion um eine konstante Funktion handelt, also um $f_1$ oder $f_2$,
oder um eine balancierte Funktion ($f_3$ oder $f_4$). Klassisch muss man die unbekannte
Funktion auf beiden Bitwerten testen, also zwei Abfragen vornehmen, 
um diese Unterscheidung treffen zu k\"onnen.

\section{Der Deutsch-Algorithmus f\"ur $N=1$}

F\"ur den einfachsten Fall $N=1$ betrachtet man
als Ausgangszustand den 2-Qubit-Zustand $|\psi_1\rangle = |0\rangle \otimes |1\rangle$.
Dieser Zustand wird in jedem seiner beiden Anteile einer Hadamard-Transformation unterworfen,
die wir in diesem Fall durch
\begin{equation}      
      H= \frac{1}{\sqrt{2}} \left( \begin{array}{cc}  1 & 1 \\  1 & -1 \end{array} \right) 
\end{equation}
darstellen. Wir erhalten somit folgenden Zustand 
\begin{equation}
    |\psi_1\rangle =  |0\rangle \otimes |1\rangle ~~ \Longrightarrow ~~  |\psi_2\rangle = 
    H |0\rangle \otimes H |1\rangle =
        \frac{1}{2} \big( |0\rangle + |1 \rangle \big) \otimes \big( |0 \rangle - |1 \rangle \big)  \, . 
\end{equation} 
Diesen Zustand kann man auch in folgender Form schreiben:
\begin{equation}
         |\psi_2\rangle = \frac{1}{2} \left( \sum_{x=0}^1 |x\rangle \right) \otimes \big(|0\rangle - |1\rangle \big) \, .
\end{equation}
Der n\"achste Schritt besteht in einer unit\"aren Implementierung des Oracles, also der
unbekannten Funktion $f$. Wir definieren f\"ur diese Implementierung $U_f$ die Transformation
\begin{equation}
             U_f | \psi_2\rangle = |\psi_3 \rangle = \frac{1}{2} \left( \sum_{x=0}^1 |x\rangle \right) \otimes
                          \big( |0 \oplus f(x)\rangle - |1 \oplus f(x) \rangle \big)  \, ,
\end{equation}
wobei $\oplus$ im Sinne einer Addition modulo 2 (bzw.\ einer XOR-Operation) zu verstehen
ist. Im Anhang \ref{sec_Deutsch_Anhang1} werden die vier
vier unit\"aren Matrizen $U_i$, die die Funktionen $f_i$ repr\"asentieren, abgeleitet:
\begin{equation}
      U_1 =  \mathbf{1} \otimes \mathbf{1} \hspace{0.8cm}  U_2= \mathbf{1} \otimes \sigma_x
      \hspace{0.8cm}  U_3= \mathrm{CNOT} \hspace{0.8cm} U_4=(\mathbf{1} \otimes \sigma_x) \mathrm{CNOT}
\end{equation}
Hierbei sind $\mathbf{1}$ die $2\times2$-Identit\"atsmatrix, 
 $\sigma_x$ die $2\times 2$ Pauli-Matrix (hier spricht man auch von einer Pauli-X-Operation),
 \begin{equation}
                \sigma_x = \left( \begin{array}{cc} 0 & 1 \\ 1 & 0 \end{array} \right) ~, 
 \end{equation}
und CNOT die $4\times 4$ unit\"are \glqq controlled NOT\grqq-Matrix, die
das zweite Bit unver\"andert l\"asst, wenn das erste Bit 0 ist, andernfalls wird das zweite Bit
invertiert:
\begin{equation}
                \mathrm{CNOT} = \left( \begin{array}{cccc} 
       1 & 0 & 0 & 0 \\ 0 & 1 & 0 & 0 \\ 0 & 0 & 0 & 1 \\ 0 & 0 & 1 & 0 \end{array} \right) ~ . 
\end{equation} 
In Matrix-Schreibweise folgt f\"ur die unit\"aren Implementationen der vier m\"oglichen Oracles:
\begin{equation}
      U_1 = \left(  \begin{array}{cccc} 
       1 & 0 & 0 & 0 \\ 0 & 1 & 0 & 0 \\ 0 & 0 & 1 & 0 \\ 0 & 0 & 0 & 1 \end{array} \right)  \hspace{0.3cm}
      U_2 = \left(  \begin{array}{cccc} 
       0 & 1 & 0 & 0 \\ 1 & 0 & 0 & 0 \\ 0 & 0 & 0 & 1 \\ 0 & 0 & 1 & 0 \end{array} \right)  \hspace{0.3cm}
      U_3 = \left(  \begin{array}{cccc} 
       1 & 0 & 0 & 0 \\ 0 & 1 & 0 & 0 \\ 0 & 0 & 0 & 1 \\ 0 & 0 & 1 & 0 \end{array} \right)  \hspace{0.3cm}
      U_4 = \left(  \begin{array}{cccc} 
       0 & 1 & 0 & 0 \\ 1 & 0 & 0 & 0 \\ 0 & 0 & 1 & 0 \\ 0 & 0 & 0 & 1 \end{array} \right)  
\end{equation}
Wann immer $f(x)=0$ ist, \"andert sich an dem Zustand nichts. Ist $f(x)=1$, werden die Rollen
von $|0\rangle$ und $|1\rangle$ im zweiten Bit vertauscht, was aber insgesamt nur ein
Minuszeichen bewirkt, da
\begin{equation}
        |0\rangle - |1\rangle ~~  \Longrightarrow ~~  |1\rangle - |0\rangle = - \big( |0\rangle - |1\rangle \big) \, .
\end{equation}
Somit k\"onnen wir f\"ur $|\psi_3\rangle$ auch schreiben:
\begin{equation}
      |\psi_3\rangle = \frac{1}{2} \left( \sum_{x=0}^1 (-1)^{f(x)} |x\rangle \right) \otimes \big( |0\rangle - |1\rangle \big) \, .  
\end{equation}
Das zweite Qubit ist nun unabh\"angig von der Funktion $f$ und wird im Folgenden nicht mehr
ben\"otigt, sodass wir es weglassen (einschlie\ss lich eines Normierungsfaktors $1/\sqrt{2}$). 
Zust\"ande, bei denen wir das letzte Qubit weglassen, kennzeichnen wir durch einen $'$, z.B.\
$|\psi_3' \rangle$. Das erste Qubit repr\"asentiert den Zustand 
\begin{equation}
               |\psi_3' \rangle =  \frac{1}{\sqrt{2}} \left( (-1)^{f(0)} |0\rangle + (-1)^{f(1)} |1\rangle \right) =
               \frac{1}{\sqrt{2}} (-1)^{f(0)} \left(  |0\rangle + (-1)^{f(0)\oplus f(1)} |1\rangle \right) \, .
\end{equation}
Hierbei wurde benutzt, dass $f(0)\oplus f(0)=0$, wir somit immer um diese Summe erweitern
k\"onnen. Nun gilt: 
\begin{equation}
             f(0)\oplus f(1)=\sigma (f) = \left\{ \begin{array}{ll} 0 & \mbox{falls $f$ konstant} \\
                   1 & \mbox{falls $f$ balanciert} \end{array} \right. 
\end{equation}
Wir erhalten
also letztendlich f\"ur das erste Qubit den Zustand
\begin{equation}
       |\psi_3'\rangle  = \left\{ \begin{array}{ll}
       {\displaystyle \frac{1}{\sqrt{2}} \big(|0\rangle + |1\rangle \big) = |+\rangle} & 
                \mbox{falls $f$ eine konstante Funktion ist} \\[0.5cm]
       {\displaystyle \frac{1}{\sqrt{2}} \big(|0\rangle - |1\rangle \big) = |- \rangle} &
       \mbox{falls $f$ balanciert ist}  \end{array} \right.  \, .
\end{equation}
Wenden wir auf diesen Zustand nochmals eine
Hadamard-Transformation an, erhalten wir f\"ur das erste Qubit den Zustand $|0\rangle$,
sofern $f$ eine konstante Funktion ist, und $|1\rangle$, sofern $f$ eine balancierte Funktion ist:
\begin{equation}
    |\psi_3' \rangle ~~  \Longrightarrow ~~ |\psi_4'\rangle = H|\psi_3'\rangle = \left\{ \begin{array}{ll}
              |0\rangle & \mbox{falls $f$ konstant} \\  |1\rangle & \mbox{falls $f$ balanciert}
              \end{array}  \right.  
\end{equation}
Wir m\"ussen also nur das erste Qubit bez\"uglich der Standardbasis $|0\rangle$ und $|1\rangle$
messen und erhalten als Ergebnis die gesuchte Antwort auf die Frage, ob es sich bei dem
Oracle um eine konstante oder balancierte Funktion handelt.

Graphisch dr\"uckt man diesen Prozess h\"aufig aus wie in Abb.\ \ref{fig_Deutsch1}.

\begin{SCfigure}[50][h!]
\begin{picture}(250,80)(-20,20)
\put(10,80){\makebox(0,0){$|0\rangle$}}
\put(10,40){\makebox(0,0){$|1\rangle$}}
\put(50,80){\makebox(0,0){$H$}}
\put(50,40){\makebox(0,0){$H$}}
\put(170,80){\makebox(0,0){$H$}}
\put(210,80){\makebox(0,0){$M$}}
\put(110,60){\makebox(0,0){\Large $U_f$}}
\put(10,20){\makebox(0,0){$|\psi_1\rangle$}}
\put(70,20){\makebox(0,0){$|\psi_2\rangle$}}
\put(150,20){\makebox(0,0){$|\psi_3\rangle$}}
\put(190,20){\makebox(0,0){$|\psi_4\rangle$}}
\put(190,60){\makebox(0,0){$|\psi_4'\rangle$}}
\thicklines
\put(20,80){\line(1,0){20}}   % obere Linie
\put(40,70){\line(1,0){20}}   % obere H-Box
\put(40,90){\line(1,0){20}}
\put(40,70){\line(0,1){20}}
\put(60,70){\line(0,1){20}}
\put(20,40){\line(1,0){20}}   % untere Linie
\put(40,30){\line(1,0){20}}   % untere H-Box
\put(40,50){\line(1,0){20}}
\put(40,30){\line(0,1){20}}
\put(60,30){\line(0,1){20}}
\put(60,80){\line(1,0){20}}   % obere Linie
\put(60,40){\line(1,0){20}}   % untere Linie
\put(80,30){\line(1,0){60}}   % gro�e Box
\put(80,90){\line(1,0){60}}
\put(80,30){\line(0,1){60}}
\put(140,30){\line(0,1){60}}
\put(140,80){\line(1,0){20}}   % obere Linie
\put(140,40){\line(1,0){80}}   % untere (leere) Linie
\put(160,70){\line(1,0){20}}   % obere H-Box
\put(160,90){\line(1,0){20}}
\put(160,70){\line(0,1){20}}
\put(180,70){\line(0,1){20}}
\put(180,80){\line(1,0){20}}   % obere Linie
\put(200,70){\line(1,0){20}}   % M-Box
\put(200,90){\line(1,0){20}}
\put(200,70){\line(0,1){20}}
\put(220,70){\line(0,1){20}}
\end{picture}
\caption{\label{fig_Deutsch1}%
Graphische Darstellung des Deutsch-Algorithmus f\"ur $N=1$. $H$ bezeichnet
Hadamard-Gatter, $U_f$ die Implementierung des Oracles $f$ und $M$ die
abschlie\ss ende Messung.}
\end{SCfigure}

\section{Der Deutsch-Jozsa-Algorithmus f\"ur $N>1$}

In diesem Abschnitt werden die Verallgemeinerungen behandelt, die bei
einem Oracle, das auf einer Folge von $N$ Qubits definiert ist, auftreten. Die unit\"are
Implementation des Oracles wird allerdings nicht explizit angegeben. Sie ist eine Verallgemeinerung
der unit\"aren Operatoren des letzten Abschnitts.

Als Ausgangszustand w\"ahlt man nun
\begin{equation}
         | \psi_1\rangle = |0\rangle \otimes |0\rangle \otimes \cdots \otimes |0\rangle \otimes |1 \rangle 
         =  |0\rangle^{\otimes N} \otimes |1\rangle  \, ,
\end{equation}
wobei $|0\rangle^{\otimes N}$ das $N$-fache Tensorprodukt des Zustands $|0\rangle$ bezeichnet. 
In einem ersten Schritt wendet man auf jedes Qubit eine Hadamard-Transformation an:
\begin{eqnarray}
   |\psi_2\rangle &=& H^{\otimes N+1} |\psi_1\rangle =
        H|0\rangle \otimes \cdots \otimes H|0\rangle \otimes H|1\rangle  \\
        &=&  \frac{1}{\sqrt{2^{N+1}}} \Big( \big( |0\rangle + |1\rangle \big) \otimes \big( |0\rangle + |1\rangle \big)
            \otimes \ldots \otimes \big( |0\rangle + |1\rangle \big) \Big) \otimes \big( |0\rangle - |1\rangle \big)  \, .
\end{eqnarray} 
Diesen Zustand kann man auch folgenderma\ss en schreiben, wenn man das $N$-fache
Tensorprodukt der $H|0\rangle$-Zust\"ande ausmultipliziert:
\begin{equation}
\label{eq_Deutsch1}
   |\psi_2\rangle = \frac{1}{\sqrt{2^{N+1}}} \left( \sum_{x=0}^{2^N-1} |x\rangle \right) 
                 \otimes \big( |0\rangle - |1\rangle \big)  \, .
\end{equation} 
Man erkennt an dieser Darstellung deutlich, dass dieser Zustand eine Superposition aus
Zust\"anden zu allen $N$-Bit-Folgen bzw.\ ihren Darstellungen durch nat\"urliche Zahlen $x$ ist. 
Man kann somit alle zu testenden F\"alle in einem einzigen Zustand kodieren. 

Die explizite unit\"are Implementierung der Oracle-Funktionen ist nun etwas komplizierter. Aber
das Ergebnis soll folgende Form haben:
\begin{equation}
        |\psi_3\rangle =   U_f |\psi_2\rangle =   \frac{1}{\sqrt{2^{N+1}}} 
        \left( \sum_{x=0}^{2^N-1} |x\rangle \right) \otimes \Big( |0\oplus f(x) \rangle - |1\oplus f(x) \rangle \Big)  \, .   
\end{equation}
Wiederum ist offensichtlich, dass sich f\"ur den Fall $f(x)=0$ an dem Zustand \"uberhaupt nichts
\"andert, und f\"ur den Fall $f(x)=1$ werden im letzten Qubit wiederum $|0\rangle$ und
$|1\rangle$ ausgetauscht, was denselben Zustand aber mit einem Minuszeichen liefert.
Technisch gesprochen ist $|x\rangle \otimes \big( |0\rangle - |1\rangle \big)$ ein Eigenzustand zu
allen $U_f$ mit dem Eigenwert $(-1)^{f(x)}$ f\"ur $x\in [0,2^N-1]$:
\begin{equation}
       U_f |x\rangle \otimes \big( |0\rangle - |1\rangle \big) 
                 = |x\rangle \otimes \big( |0\oplus f(x)\rangle - |1\oplus f(x) \rangle \big) 
                     = (-1)^{f(x)} |x \rangle \otimes \big( |0\rangle - |1\rangle \big) \, .
\end{equation} 
Damit k\"onnen wir $|\psi_3 \rangle$ auch in folgender Form schreiben:
\begin{equation}
        | \psi_3 \rangle =   \frac{1}{\sqrt{2^{N+1}}} 
        \left( \sum_{x=0}^{2^N-1} (-1)^{f(x)}  |x\rangle \right) \otimes \big( |0 \rangle - |1  \rangle \big)  \, .
\end{equation}
Das letzte Qubit ist von dem Oracle unabh\"angig und kann weggelassen werden.
F\"ur die ersten $N$ Qubits gilt: Ist $f(x)=f$ eine konstante Funktion, \"andert sich der Zustand nicht,
der Zustandsvektor wird lediglich mit $(-1)^f$ multipliziert:
\begin{equation}
       |\psi_3' \rangle = \frac{1}{\sqrt{2^{N}}} (-1)^{f}  \left( \sum_{x=0}^{2^N-1}  |x\rangle \right) 
       \hspace{0.8cm} \mbox{($f$ konstant)}\, .
\end{equation}
Durch eine abschlie\ss ende Hadamard-Transformation auf jedem Qubit wird aus diesem
Zustand bis auf den Faktor $(-1)^f$ wieder der Ausgangszustand:
\begin{equation}
      | \psi_4' \rangle = H^{\otimes N} |\psi_3'\rangle =  (-1)^f | 0 \rangle^{\otimes N}
\end{equation}
Wir zeigen nun allgemeiner die folgende Aussage:
\begin{equation}
          \langle \psi_4' | 0 \rangle^{\otimes N}  = \left\{ \begin{array}{ll}
          (-1)^f  & \mbox{falls $f$ konstant} \\
          0 & \mbox{falls $f$ balanciert} \end{array} \right.  \, .
\end{equation}
Daraus folgt: Die Wahrscheinlichkeit $w(|0\rangle^{\otimes N})$, am Ausgang unseres
Quantencomputers den Zustand $|0\rangle^{\otimes N}$ zu messen -- alle $N$ Messger\"ate an
den jeweiligen Ausg\"angen messen den Zustand $|0\rangle$ -- ist gegeben durch:
\begin{equation}
  w(|0\rangle^{\otimes N}) = \left\{  \begin{array}{ll}
            1 & \mbox{falls $f$ konstant} \\ 0 & \mbox{falls $f$ balanciert} \end{array} \right. \, .
\end{equation}
Zum Beweis betrachten wir das Skalarprodukt von $|\psi_4'\rangle = H^{\otimes N}|\psi_3' \rangle$
mit dem Zustand $|0\rangle^{\otimes N}$, wenden die Hadamard-Transformation nun
aber auf $|0\rangle^{\otimes N}$ an. Dies liefert wie schon in Gl.\ \ref{eq_Deutsch1} eine
Summe \"uber alle nat\"urlichen Zahlen von 0 bis $2^N-1$ im Bin\"arcode:
\begin{equation}
        \langle \psi_4' | 0\rangle^{\otimes N}  = \frac{1}{\sqrt{2^N}}
      \sum_{x=0}^{2^N-1} (-1)^{f(x)} \langle x |H^{\otimes N}| 0 \rangle^{\otimes N}  
      = \frac{1}{2^N} \sum_{x,y=0}^{2^N-1} (-1)^{f(x)} \langle x |y \rangle \, .
\end{equation}
Da Zust\"ande zu verschiedenen Bitfolgen orthogonal sind, gilt
$\langle x|y\rangle = \delta_{xy}$, wir erhalten also nur einen Beitrag f\"ur $x=y$:
\begin{equation}
      \langle \psi_4' | 0\rangle^{\otimes N}  = \frac{1}{2^N}
      \sum_{x=0}^{2^N-1} (-1)^{f(x)} \, .
\end{equation}
Ist $f$ konstant, k\"onnen wir den Faktor $(-1)^f$ vor die Summe ziehen und erhalten als
Ergebnis 1. Ist $f$ balanciert, erhalten wir eine Summe von $+1$ und $-1$, wobei beide
Zahlen gleich h\"aufig auftreten. Insgesamt ist in diesem Fall das Ergebnis daher 0.
Damit ist unsere obige Behauptung bewiesen. 
Zur graphischen Darstellung des Deutsch-Jozsa-Algorithmus f\"ur beliebige $N$, 
siehe Abb.\ \ref{fig_Deutsch2}.

\begin{SCfigure}[50][h!]
\begin{picture}(250,120)(-10,20)
\put(10,120){\makebox(0,0){$|0\rangle$}}
\put(10,80){\makebox(0,0){$|0\rangle$}}
\put(10,40){\makebox(0,0){$|1\rangle$}}
\put(50,120){\makebox(0,0){$H$}}
\put(50,80){\makebox(0,0){$H$}}
\put(50,40){\makebox(0,0){$H$}}
\put(170,120){\makebox(0,0){$H$}}
\put(220,120){\makebox(0,0){$M$}}
\put(170,80){\makebox(0,0){$H$}}
\put(220,80){\makebox(0,0){$M$}}
\put(110,80){\makebox(0,0){\Large $U_f$}}
\put(10,20){\makebox(0,0){$|\psi_1\rangle$}}
\put(70,20){\makebox(0,0){$|\psi_2\rangle$}}
\put(150,20){\makebox(0,0){$|\psi_3\rangle$}}
\put(195,20){\makebox(0,0){$|\psi_4\rangle$}}
\put(195,65){\makebox(0,0){$|\psi_4'\rangle$}}
\put(5,103){\makebox(0,0){$\vdots$}}
\put(50,103){\makebox(0,0){$\vdots$}}
\put(165,103){\makebox(0,0){$\vdots$}}
\put(220,103){\makebox(0,0){$\vdots$}}
\put(15,100){\makebox(0,0){$\otimes N$}}
\put(175,100){\makebox(0,0){$\otimes N$}}
\thicklines
\put(20,120){\line(1,0){20}}   % obere Linie
\put(40,110){\line(1,0){20}}   % obere H-Box
\put(40,130){\line(1,0){20}}
\put(40,110){\line(0,1){20}}
\put(60,110){\line(0,1){20}}
\put(20,80){\line(1,0){20}}   % mittlere Linie
\put(40,70){\line(1,0){20}}   % mittlere H-Box
\put(40,90){\line(1,0){20}}
\put(40,70){\line(0,1){20}}
\put(60,70){\line(0,1){20}}
\put(20,40){\line(1,0){20}}   % untere Linie
\put(40,30){\line(1,0){20}}   % untere H-Box
\put(40,50){\line(1,0){20}}
\put(40,30){\line(0,1){20}}
\put(60,30){\line(0,1){20}}
\put(60,120){\line(1,0){20}}   % obere Linie
\put(60,80){\line(1,0){20}}   % mittlere Linie
\put(60,40){\line(1,0){20}}   % untere Linie
\put(80,30){\line(1,0){60}}   % gro�e Box
\put(80,130){\line(1,0){60}}
\put(80,30){\line(0,1){100}}
\put(140,30){\line(0,1){100}}
\put(140,120){\line(1,0){20}}   % obere Linie
\put(140,40){\line(1,0){90}}   % untere (leere) Linie
\put(160,110){\line(1,0){20}}   % obere H-Box
\put(160,130){\line(1,0){20}}
\put(160,110){\line(0,1){20}}
\put(180,110){\line(0,1){20}}
\put(180,120){\line(1,0){30}}   % obere Linie
\put(210,110){\line(1,0){20}}   % M-Box
\put(210,130){\line(1,0){20}}
\put(210,110){\line(0,1){20}}
\put(230,110){\line(0,1){20}}

\put(140,80){\line(1,0){20}}   % obere Linie
\put(160,70){\line(1,0){20}}   % obere H-Box
\put(160,90){\line(1,0){20}}
\put(160,70){\line(0,1){20}}
\put(180,70){\line(0,1){20}}
\put(180,80){\line(1,0){30}}   % obere Linie
\put(210,70){\line(1,0){20}}   % M-Box
\put(210,90){\line(1,0){20}}
\put(210,70){\line(0,1){20}}
\put(230,70){\line(0,1){20}}
\end{picture}
\caption{\label{fig_Deutsch2}%
Graphische Darstellung des Deutsch-Algorithmus f\"ur $N>1$. Die oberen $N$ Qubits
kodieren nach dem Hadamard-Gatter s\"amtliche $N$-Bit-Folgen. $U_f$ bezeichnet wieder
die Implementierung des Oracles und $M$ die abschlie\ss enden Messungen an jedem
Qubit.}
\end{SCfigure}

Man kann daher mit diesem Algorithmus
mit einer einzige Abfrage des Oracles die Information erhalten, ob die unbekannte Funktion
konstant oder balanciert ist. 

%Man beachte, dass bei keinem Schritt dieses Quantenalgorithmus' ein verschr\"ankter
%Zustand auftritt. 

\section{Das Penny-Flip-Spiel}

Das Penny-Flip-Spiel von David Meyer ist eine vereinfachte Version des Deutsch-Algorithmus 
f\"ur $N=1$, bei dem die Quantenversion eine Gewinnstrategie bei einem einfachen Spiel
erm\"oglicht, die klassisch nicht m\"oglich ist \cite{Meyer,Greinert}.

Die klassischen Spielregeln sind folgende: Alice darf eine M\"unze auf den Tisch legen,
so dass \glqq Kopf\grqq\ oder \glqq Zahl\grqq\ oben liegen. Anschlie\ss end darf Bob die 
M\"unze von Alice entweder umdrehen oder so liegen lassen, wie Alice sie hingelegt hat. Alice sieht
dabei nicht, was Bob getan hat. Alice darf nun die M\"unze (ohne diese anzuschauen oder
abzutasten, d.h.,
sie wei\ss\ nicht, ob Bob die M\"unze gedreht hat oder nicht) nochmals drehen oder auch
belassen, wie sie ist. Alice hat gewonnen, wenn die M\"unze abschlie\ss end \glqq Kopf\grqq\
zeigt, Bob hat gewonnen, wenn die M\"unze \glqq Zahl\grqq\ zeigt. Gibt es f\"ur Alice
eine Gewinnstrategie?

Offensichtlich gibt es bei dieser klassischen Version des Spiels f\"ur Alice keine
Gewinnstrategie, da sie nicht wissen kann, ob Bob die M\"unze gedreht hat oder nicht.
In einer Quantenversion dieses Spiels hat Alice jedoch eine 100\%-ige Gewinnchance.
In der Quantenversion kann Alice die M\"unze in einem Superpositionszustand aus
\glqq Kopf\grqq\ und \glqq Zahl\grqq\ pr\"aparieren. Sie kann beispielsweise die M\"unze
in dem Zustand
\begin{equation}
         |\psi_2\rangle = H | \mbox{Kopf} \rangle = \frac{1}{\sqrt{2}} ( |\mbox{Kopf} + | \mbox{Zahl}\rangle )
\end{equation}
pr\"aparieren, wobei $H$ wieder ein Hadamard-Gatter f\"ur den Zustand der M\"unze
darstellt.

Bob darf nun eine von zwei Transformationen durchf\"uhren: Entweder er l\"asst die M\"unze
unver\"andert, also 
\begin{equation}
     | \mbox{Zahl} \rangle ~~\Longrightarrow ~~| \mbox{Zahl} \rangle 
      \hspace{0.8cm}  \mbox{und} \hspace{0.8cm} 
           | \mbox{Kopf} \rangle ~~\Longrightarrow ~~| \mbox{Kopf} \rangle  \, ,
\end{equation}
oder er dreht die M\"unze:
\begin{equation}            
     | \mbox{Zahl} \rangle ~~\Longrightarrow ~~| \mbox{Kopf} \rangle 
      \hspace{0.8cm}  \mbox{und} \hspace{0.8cm} 
           | \mbox{Kopf} \rangle ~~\Longrightarrow ~~| \mbox{Zahl} \rangle \, .
\end{equation}
Hier wird ausgenutzt, dass der von Alice pr\"aparierte Superpositionszustand ein
Eigenzustand sowohl zur Identit\"atstransformation als auch zur Vertauschungsoperation
(Pauli-X) ist:
\begin{equation}
     \left( \begin{array}{cc} 1 & 0 \\ 0 & 1 \end{array} \right)
     \frac{1}{\sqrt{2}}  \left( \begin{array}{c} 1 \\ 1 \end{array} \right)  =
          \frac{1}{\sqrt{2}}  \left( \begin{array}{c} 1 \\ 1 \end{array} \right)  \hspace{1.5cm}
     \left( \begin{array}{cc} 0 & 1 \\ 1 & 0 \end{array} \right)
     \frac{1}{\sqrt{2}}  \left( \begin{array}{c} 1 \\ 1 \end{array} \right)  =
          \frac{1}{\sqrt{2}}  \left( \begin{array}{c} 1 \\ 1 \end{array} \right)  ~~. 
\end{equation}
Egal f\"ur welche der beiden Operationen sich Bob entscheidet,
der Zustand wird durch seine Operation nicht ver\"andert. Abschlie\ss end
kann Alice nochmals eine Hadamard-Transformation auf den Zustand anwenden und
erh\"alt mit Sicherheit den Zustand $|\mbox{Kopf}\rangle$, unabh\"angig davon, was Bob
gemacht hat. Insofern hat Alice immer gewonnen.

\section{Anhang}

\subsection{Die unit\"are Implementation des Oracles}
\label{sec_Deutsch_Anhang1}

In diesem Anhang soll gezeigt werden, wie man f\"ur den Fall $N=1$ zu den Oracle-Funktionen die
unit\"aren Matrizen konstruiert und wie diese durch Quantengatter realisiert werden
k\"onnen.

Die Funktion $f$ soll durch folgende Vorschrift realisiert werden:
\begin{equation}
         f : | x \rangle \otimes |y \rangle ~ \longrightarrow ~ |x \rangle \otimes |y \oplus f(x) \rangle \, ,
\end{equation}
wobei $x,y\in \{0,1\}$ und die $\oplus$-Operation als \glqq Addition modulo 2\grqq\ oder
XOR-Operation zu verstehen sind. Als 4-Vektor wird $|x\rangle \otimes |y\rangle$ folgenderma\ss en
repr\"asentiert:
\begin{equation}
    |0\rangle \otimes |0\rangle = \left( \begin{array}{c} 1 \\ 0 \\ 0 \\ 0 \end{array} \right) \hspace{0.5cm}
    |0\rangle \otimes |1\rangle = \left( \begin{array}{c} 0 \\ 1 \\ 0 \\ 0 \end{array} \right) \hspace{0.5cm}
    |1\rangle \otimes |0\rangle = \left( \begin{array}{c} 0 \\ 0 \\ 1 \\ 0 \end{array} \right) \hspace{0.5cm}
    |1\rangle \otimes |1\rangle = \left( \begin{array}{c} 0 \\ 0 \\ 0 \\ 1 \end{array} \right)     
\end{equation}
Dies folgt den allgemeinen Konventionen:
\begin{equation}
   |0\rangle = \left( \begin{array}{c} 1 \\ 0 \end{array} \right) \hspace{0.5cm}
   |1\rangle = \left( \begin{array}{c} 0 \\ 1 \end{array} \right) \hspace{0.5cm} \mbox{und} \hspace{0.5cm}
   \left( \begin{array}{c} x_1 \\ x_2 \end{array} \right) \otimes   \left( \begin{array}{c} y_1 \\ y_2 \end{array} \right)
        =  \left( \begin{array}{c} x_1 y_1 \\ x_1 y_2 \\ x_2 y_1 \\ x_2 y_2 \end{array} \right) \, . 
\end{equation}
Wir konstruieren nun die Darstellungen der vier Oracles zu $N=1$. Man beachte dabei, dass
sich das erste Qubit nie \"andert, sondern nur das zweite Qubit -- eventuell in Abh\"angigkeit von dem
Zustand des ersten Qubits. Der Einfachheit schreibe ich im Folgenden $|x,y\rangle$ statt
$|x\rangle \otimes |y \rangle$. 

\subsubsection{$f_1$}

Die Funktion $f_1$ hat immer
den Wert $f_1(x)=0$ und \"andert den Zustand daher nicht. Sie wird durch die $4\times 4$-Identit\"atsmatrix
repr\"asentiert:
\begin{equation}
         U_1 = 
          \left(  \begin{array}{cccc} 1 & 0 & 0 & 0 \\ 0 & 1 & 0 & 0 \\ 0 & 0 & 1 & 0 \\ 0 & 0 & 0 & 1 \end{array} \right)
           = {\bf 1} \otimes {\bf 1} \, .
\end{equation}

\subsubsection{$f_2$}

Die Funktion $f_2$ hat immer den Wert $f_2(x)=1$ und tauscht daher im zweiten Qubit die
Werte von $0$ und $1$ aus. Das bedeutet:
\begin{equation}
    |0,0\rangle \mapsto |0,1\rangle \hspace{0.8cm}  |0,1\rangle \mapsto |0,0\rangle \hspace{0.8cm}
    |1,0\rangle \mapsto |1,1\rangle \hspace{0.8cm}  |1,1\rangle \mapsto |0,1\rangle \, .       
\end{equation} 
Die Matrix zu dieser Abbildung ist
\begin{equation}
      U_2 = \left(  \begin{array}{cccc} 
       0 & 1 & 0 & 0 \\ 1 & 0 & 0 & 0 \\ 0 & 0 & 0 & 1 \\ 0 & 0 & 1 & 0 \end{array} \right) 
          = {\bf 1} \otimes \sigma_x   \, . 
\end{equation}

\subsubsection{$f_3$}

Die Abbildung $f_3(x)=x$ l\"asst das zweite Qubit unver\"andert, wenn $x=0$ ist, und ersetzt
es durch seine Negation, wenn $x=1$ ist. Das bedeutet f\"ur die Standardbasisvektoren:
\begin{equation}
    |0,0\rangle \mapsto |0,0\rangle \hspace{0.8cm}  |0,1\rangle \mapsto |0,1\rangle \hspace{0.8cm}
    |1,0\rangle \mapsto |1,1\rangle \hspace{0.8cm}  |1,1\rangle \mapsto |1,0\rangle \, .       
\end{equation} 
Diese Transformation wird durch folgende Matrix repr\"asentiert:
\begin{equation}
      U_3 = \left(  \begin{array}{cccc} 
       1 & 0 & 0 & 0 \\ 0 & 1 & 0 & 0 \\ 0 & 0 & 0 & 1 \\ 0 & 0 & 1 & 0 \end{array} \right) = \mbox{CNOT} \, .
\end{equation}
Man bezeichnet diese Matrix als \glqq Controlled NOT\grqq, da die NOT-Operation auf dem
zweiten Qubit durch das erste Qubit kontrolliert wird.

\subsection{$f_4$} 

Die vierte Abbildung $f_4(x)= \neg x$ \"andert das zweite Qubit, sofern das erste sich im
Zustand $|0\rangle$ befindet. Es handelt sich daher bei der Implementierung ebenfalls 
um ein CNOT-Gatter, allerdings ist der Control-Parameter nun vertauscht. F\"ur die Basis
bedeutet dies: 
\begin{equation}
    |0,0\rangle \mapsto |0,1\rangle \hspace{0.8cm}  |0,1\rangle \mapsto |0,0\rangle \hspace{0.8cm}
    |1,0\rangle \mapsto |1,0\rangle \hspace{0.8cm}  |1,1\rangle \mapsto |1,1\rangle \, .       
\end{equation} 
Die zugeh\"orige Transformationsmatrix ist:
\begin{equation}
      U_4 = \left(  \begin{array}{cccc} 
       0 & 1 & 0 & 0 \\ 1 & 0 & 0 & 0 \\ 0 & 0 & 1 & 0 \\ 0 & 0 & 0 & 1 \end{array} \right) = 
       ({\bf 1} \otimes \sigma_3) \mbox{CNOT} \, .
\end{equation}

Alle vier Oracles lassen sich somit durch geeignete Hintereinanderschaltung der beiden
Gatter Pauli-X und CNOT darstellen.


\begin{thebibliography}{99}
\bibitem{Deutsch} Deutsch, David (1985); \textit{Quantum Theory, the Church-Turing Principle and the
       Universal Quantum Computer}; Proceedings of the Royal Society of London A.\ 400 (1818):
       97--117.
\bibitem{DJ} Deutsch, David, Jozsa, Richard (1992); \textit{Rapid solutions of problems by quantum
         computation}. Proceedings of the Royal Society of London A.\ 439 (1907): 553--558.
\bibitem{Ekert} Cleve, R., Ekert, A, Macchiavello, C., Mosca, M, (1998); \textit{Quantum algorithms revisited};
       Proceedings of the Royal Society of London A.\ 454 (1969): 339--354.
\bibitem{Meyer} Meyer, David A. (1999); \textit{Quantum Strategies}; Phys.\ Rev.\ Lett.\ 82, 1052       
\bibitem{Greinert} M\"uller, Rainer; Greinert, Franziska (2023); \textit{Quantum Penny Flip -- Spielen
         mit dem Quantencomputer}; plusLucis 2/2023 Quantenphysik. \textit{Playing with a Quantum
         Computer}, arXiv: 2108.06271.                           
\end{thebibliography}

\end{document}

